\chapter{Conclusions and Future Work}

\section{Final Answer to the Project Question}

The project started from a precise research question:

\begin{quote}
    Can a feed-forward neural network accurately approximate the Black-Scholes pricing function for European call options?
\end{quote}

The answer is positive within the controlled setting studied in the report.
The final model learns the map

\begin{equation}
    (S_0, K, T, r, \sigma, S_0/K)
    \longmapsto
    C_{BS}
\end{equation}

with high accuracy on held-out synthetic observations.
Using 100,000 synthetic contracts, the final neural network achieves:

\begin{equation}
    \mathrm{MAE}=0.04290,
    \qquad
    \mathrm{RMSE}=0.06208,
    \qquad
    R^2=0.999993.
\end{equation}

These values show that the neural network reproduces the analytical Black-Scholes price very closely over the sampled parameter domain.
The result is also supported visually by the true-vs-predicted plot and by the error diagnostics discussed in Chapter~7.

\begin{quote}
    The project demonstrates that a supervised feed-forward neural network can act as an accurate pricing surrogate for the Black-Scholes pricing function in a controlled synthetic setting.
\end{quote}

This conclusion is intentionally limited to the scope of the experiment.
The model is not claimed to price real exchange-traded options, and it is not presented as a replacement for the closed-form Black-Scholes formula.
Instead, the contribution is to show how a neural network can learn a known pricing map and provide fast inference once trained.

\section{Main Findings}

For compactness, the main findings are collected in Table~\ref{tab:final-findings} before discussing their interpretation.

\begin{table}[H]
    \centering
    \begin{tabular}{
        >{\raggedright\arraybackslash}p{0.31\linewidth}
        >{\raggedright\arraybackslash}p{0.31\linewidth}
        >{\raggedright\arraybackslash}p{0.28\linewidth}
    }
        \toprule
        Finding & Evidence & Interpretation \\
        \midrule
        The neural network approximates Black-Scholes accurately
        & low MAE and RMSE, $R^2$ close to one
        & the MLP learns the pricing map on the sampled domain \\
        \projecttablerowsep
        Moneyness and SiLU improve the final model
        & intermediate experiments and improved final run
        & financial feature engineering and smooth activations help regression \\
        \projecttablerowsep
        Neural inference is faster than Monte Carlo after training
        & runtime benchmark on 1,000 pricing queries
        & the trained model is useful as a repeated-evaluation surrogate \\
        \projecttablerowsep
        Noisy labels reduce accuracy but do not destroy the approximation
        & controlled noisy-target experiments
        & the model shows moderate robustness, within a synthetic noise setup \\
        \projecttablerowsep
        SVR is a useful classical reference
        & reduced-scale SVR benchmark
        & classical ML performs well, but is not the scalable main pipeline \\
        \bottomrule
    \end{tabular}
    \caption{Summary of the main project findings.}
    \label{tab:final-findings}
\end{table}

The first finding is the core result of the project.
It answers the research question directly.
The remaining findings explain why the final model was selected, how it compares with classical methods, and how robust the approach is when the training labels are perturbed.

\section{Role of the Neural Network}

The neural network has a specific role in the project.
It is not the source of the theoretical price.
The theoretical price is still given by the Black-Scholes formula:

\begin{equation}
    C_{BS}
    =
    S_0 N(d_1)
    -
    K e^{-rT} N(d_2).
\end{equation}

The neural network is instead a learned approximation:

\begin{equation}
    \hat{C}_{NN}
    =
    f_\theta(x).
\end{equation}

This distinction matters.
When the analytical formula is available, Black-Scholes remains exact and faster.
The neural network becomes interesting as a surrogate-model prototype: it shows how a learned model can approximate a pricing map and then produce repeated predictions through fast inference.

The Monte Carlo comparison clarifies this role.
Monte Carlo remains a flexible numerical method, especially in settings where closed-form solutions are unavailable.
However, it requires repeated sampling.
The neural network requires training first, but after training it produces deterministic predictions without simulating paths.

\section{Limitations}

The main limitations of the project are not accidental weaknesses.
They define the boundary of the research question.

\begin{itemize}
    \item The dataset is synthetic and generated under Black-Scholes assumptions.
    \item Only European call options are considered.
    \item The target pricing function is known analytically.
    \item The input domain is bounded by the sampled parameter ranges.
    \item No-arbitrage constraints beyond non-negativity are not explicitly enforced.
    \item Noisy targets are controlled perturbations and should not be interpreted as real market prices.
    \item Real exchange-traded option data are not used, because that would change the research question from function approximation to market-price modeling.
\end{itemize}

The most important limitation is the synthetic setting.
This project answers a clean mathematical question:

\begin{equation}
    \text{Can a neural network approximate } C_{BS}?
\end{equation}

It does not answer the different empirical-market question:

\begin{equation}
    \text{Can a neural network explain observed option market prices?}
\end{equation}

The second question would require additional variables, careful data cleaning, treatment of bid-ask spreads, dividends, liquidity effects, and implied-volatility structure.
For this reason, keeping the current project synthetic is methodologically consistent.

\section{Future Work}

Several extensions are possible and they can be grouped according to how much they change the current research question:

\begin{itemize}
    \item Additional neural architectures, such as deeper multilayer perceptrons or residual networks, while keeping the same Black-Scholes approximation task.
    \item More detailed regime-based error analysis, separating deep out-of-the-money, at-the-money, and deep in-the-money contracts.
    \item Greeks estimation through automatic differentiation, extending the surrogate from prices to risk sensitivities.
    \item Heston or other stochastic-volatility models, which would change the underlying pricing dynamics.
    \item Basket or path-dependent options, where the dimensionality and computational relevance of surrogate pricing become more pronounced.
    \item American options, which introduce early-exercise features and require a different pricing framework.
    \item Real exchange-traded option data, which would change the project from approximating Black-Scholes to modeling market prices.
\end{itemize}

The most natural first extension would be a more detailed regime-based error analysis.
For example, one could separately study deep out-of-the-money, at-the-money, and deep in-the-money contracts.
This would keep the current research question unchanged while improving the understanding of where the neural surrogate is most reliable.

A second extension would be Greeks estimation.
Since the neural network is differentiable, sensitivities could be estimated by automatic differentiation:

\begin{equation}
    \Delta_{NN}
    =
    \frac{\partial f_\theta}{\partial S_0},
    \qquad
    \mathcal{V}_{NN}
    =
    \frac{\partial f_\theta}{\partial \sigma}.
\end{equation}

These quantities could then be compared with analytical Black-Scholes Greeks.
This would be a natural continuation because option pricing is often used together with risk sensitivities.

More ambitious extensions include stochastic volatility, basket options, American options, and real market data.
These are valuable directions, but they would move beyond the initial purpose of this project.
They should therefore be treated as new research questions rather than small modifications of the current experiment.

\section{Final Remarks}

The project combines stochastic modeling, option pricing, Monte Carlo simulation, supervised learning, and reproducible software engineering.
The final result is a complete experimental pipeline: synthetic data generation, analytical pricing, Monte Carlo benchmarking, neural-network training, SVR comparison, noisy-target robustness, runtime analysis, plots, tests, documentation, and a LaTeX report.

The main lesson is that neural networks can be useful in quantitative finance when their role is clearly defined.
In this project, the neural network is not a black-box trading model.
It is a controlled function approximator trained to reproduce a known pricing rule.
This makes the experiment mathematically interpretable, computationally reproducible, and suitable as a foundation for more complex pricing problems.
